%% Discussion (MedIA revision, v2 pipeline). Numbers are those of Section 5 (sections/results.tex);
%% the case studies (sections/case_study.tex) use v1 numbers, labeled as such.
\section{Discussion}
\label{sec:discussion}

\subsection{What a domain-complete, seed-replicated evaluation shows}
\label{sec:disc-summary}

Four findings stand out.

First, a held-out score is mostly a statement about the held-out domain. For
most arms, the domain explains most of the variance of a single run, and the
same method can score near chance at one scanner and above 0.9 at another
(Section~\ref{sec:results-domains}). The worst domain also differs between
methods. A result on one official test split therefore describes one domain,
and a different split could reverse a comparison.

Second, with three to seven domains, most differences between methods cannot
be resolved statistically. No standard-tier arm beats ERM after correcting for
multiple comparisons, even when its mean gain is 0.145 AUROC and it improves
every domain. On the tumor cohorts, gains are large but vary strongly between
domains, so between 7 and 53 domains would be needed to confirm a gain of 0.05. On the
mitosis cohorts, differences are measured precisely but are small. Adding seeds
does not fix this, because the domain-by-method interaction dominates. This
agrees with DomainBed, where no algorithm beat ERM by a margin the authors
found significant \citep{gulrajani2021domainbed}, and with the large ERM seed
variance reported for Camelyon17 in WILDS \citep{koh2021wilds}.

Third, rankings are uncertain within a cohort and change between cohorts. At
the standard tier they agree more between cohorts that share a task than
between cohorts that share a shift type. This reverses the thesis of our own
earlier version, which held that the type of shift decides which method works
best. That thesis rested on two cohorts that differed in task as well as in
shift type, and the factorial design was built to separate the two. With four
cohorts, one per cell, the design can describe the pattern but cannot prove it,
and the compact tier does not clearly show it. We read the result as a warning,
not a law. A method's rank on one task says little about its rank on another,
even when the shift looks similar.

Fourth, implementation and tuning move results as much as the choice of
method. Estimating the stain matrix per tile instead of per slide turned the
fourth-best compact-tier arm on Camelyon17 into the worst. Inheriting a learning rate tuned
for another initialization made Barlow Twins fine-tuning fail on MIDOG++. The
defects that we found in our own earlier pipeline (Section~\ref{sec:casestudy})
are of the same kind. A comparison of methods is only as good as the weakest
implementation in it.

The method-level results are consistent with earlier work. Under the primary selection rule, none of the six
domain-generalization objectives ranked in the top five on any cohort. This is
consistent with DomainBed and WILDS, where such objectives did not reliably
beat ERM \citep{gulrajani2021domainbed,koh2021wilds}. Color
augmentation was consistently good, as \citet{tellez2019stain} found for HSV
and HED augmentation. Pathology foundation models were strong on the tumor
cohorts, but their frozen features were weak for mitotic figures, and two of
the four Lunit backbones were weaker frozen encoders than ImageNet ConvNeXt-T.
Domain-aligned pre-training beats ImageNet pre-training on average
\citep{kang2023benchmarking}, but which method helps varies across datasets and
shifts \citep{wiles2022fine}, and in our data it also varied with the task.

\subsection{A roadmap for models that work across sites without retraining}
\label{sec:roadmap}

Every number in this study answers a practical question: how well does a
model trained once on some sites work at a site it has never seen? No model
was trained on labeled data from its test site, and apart from the test-time
adaptation arms, no model was changed at all. The study therefore
doubles as a test of the ``train once, deploy at many sites'' strategy that
practitioners want. At the standard scale, the best method of each cohort
reached a worst-site AUROC of 0.884 to 0.986 without any retraining, against
0.534 to 0.895 for ERM. No single method was best everywhere, and no
improvement over ERM could be confirmed statistically with the sites
available. A practitioner can therefore improve the odds that a model
transfers, and can find out whether it does, but cannot assume it.

Figure~\ref{fig:roadmap} and Table~\ref{tab:roadmap} turn the results into six
steps. Each step carries an evidence grade. \emph{Strong} means that the
finding holds in all four cohorts at the standard tier. \emph{Moderate} means a
consistent direction that is not significant after correction, or a result
that holds with exceptions. \emph{Mixed}
means large gains in some cohorts and none in others. Step~6 uses the site acceptance test proposed in
Section~\ref{sec:sat}.

\begin{figure}[pos=tbp]
  \centering
  \includegraphics[width=\linewidth]{fig_roadmap}
  \caption{A six-step roadmap for pathology models that should work at new
  sites without retraining, with the evidence for each step from this study
  (standard tier, ID-validation selection) and its evidence grade. Details are
  in Table~\ref{tab:roadmap}.}
  \label{fig:roadmap}
\end{figure}

\paragraph{Step 1: define the target before training} Decide which task,
which sites and which scanners the model is meant for, and fix the minimum
performance that every site must reach. Judge models by their worst site.
Averages hid a near-chance scanner on the canine cohort (mean AUROC 0.785,
worst scanner 0.534 for ERM), and the hardest site differed between methods.

\paragraph{Step 2: collect sites, not only slides} Gather development data
from as many sites and scanners as possible, and split it by site. Most of the
variance of a held-out score comes from the site (median ICC 0.72--0.95), so
extra slides or extra training runs from the same sites do little. A pilot
study gives $\sigma_\Delta$, and Eq.~\eqref{eq:power} then says how many sites
a validation needs. In our tumor cohorts, confirming a gain of 0.05 AUROC
needed 7 to 53 sites.

\paragraph{Step 3: train once, with a strong recipe} Start from a pathology
foundation model, fine-tune it end to end with H\&E color jitter, and tune the
learning rate for that initialization. In our data, color jitter improved mean
AUROC over ERM in 18 of 20 held-out domains at the standard tier, and the
fine-tuned DINO ViT with jitter had the best worst-site AUROC averaged over
the four cohorts (0.895). Two warnings apply. Frozen foundation-model features
were not enough for mitotic figures (the frozen probes ranked 18th to 21st on
both mitosis cohorts), so frozen features must be checked on the task at hand.
And an initialization with the wrong learning rate can fail outright: Barlow
Twins fine-tuning with the ImageNet learning rate reached 0.715 mean AUROC on
MIDOG++, against 0.891 with its own rate. Domain-generalization objectives
deserve low priority; under the primary selection rule, none of the six
reached the top five on any cohort.

\paragraph{Step 4: adapt to the new site without labels, if the pipeline
allows it} Two techniques need neither labels nor any change to the trained
weights. They use only unlabeled images of the new site at prediction time.
Re-estimating the batch-normalization statistics (BN-adapt) helped most where
the shift was severe. At the standard tier it raised mean AUROC by 0.120 on the
canine scanners, at all five of them, and by 0.015 and 0.022 on Camelyon17 and
MIDOG~2021, but it did nothing on MIDOG++ ($-0.004$). Per-slide stain
normalization gave smaller gains than color jitter in three of the four
cohorts at the standard tier. Both
techniques must be implemented with care. The test images must arrive in
shuffled, mixed batches, because a batch of one class corrupts the statistics
(Section~\ref{sec:setup-arms}). Stain must be estimated from enough tissue: a
per-tile estimate failed at a whole Camelyon17 center (0.557 against 0.976 for
the per-slide estimate; Supplementary Section~\ref{S-app:macenko}). TENT, by contrast,
updates model parameters at test time and is thus a form of fine-tuning.

\paragraph{Step 5: select and validate without touching the target sites}
Choose hyperparameters and checkpoints on held-out data from the training
sites only. Validate by holding out whole sites in turn with at least three
seeds per model. Report every site, the worst site, the seed spread, the
calibration at each site and confidence intervals that account for the
training data shared between folds. At the standard scale, the choice of
selection rule barely changed the rankings ($\tau_b\ge0.83$). With the small
compact model it did ($\tau_b$ down to 0.32), which is one more reason not to
draw conclusions from toy models. Calibration also varies: on Camelyon17 alone,
the mean ECE of standard-tier methods ranged from 0.027 to 0.186.

\paragraph{Step 6: accept each new site with a small labeled sample, then
monitor} Before a model is used at a new site, label a small random sample of
that site's patches, spread over many slides, and run the site acceptance test
of Section~\ref{sec:sat} against the target of Step~1. The development results
of Step~5 supply its prior, and its planner gives the number of labels to
collect. At the standard tier, the test kept false acceptance at or below 2.1\%
and, for targets of 0.80 and 0.85, needed about half as many labels as local
labels alone. The development
results alone were not enough: they would have accepted 8.6\% of the sites
that fell below a target of 0.80. Where the development sites score far above the
target, a site that fails is a surprise to the prior, so at least 100 labels
should be collected; the robust variant, which we added after seeing this
case, lowers the risk further. Where they score near or below the target, the
prior adds little and local labels decide. After acceptance, inputs and outputs
should be monitored; we did not test monitoring.

\begin{table}[width=\linewidth,pos=tbp]
\caption{The roadmap of Fig.~\ref{fig:roadmap} with its evidence. Numbers are
standard-tier results under ID-validation selection unless stated. Grades:
strong, holds in all four cohorts at the standard tier; moderate, consistent
direction but not significant after correction, or holding with exceptions;
mixed, large gains in some cohorts and none in others.}
\label{tab:roadmap}
\footnotesize
\begin{tabular*}{\tblwidth}{@{\extracolsep{\fill}}>{\raggedright\arraybackslash}p{2.1cm}>{\raggedright\arraybackslash}p{5.2cm}>{\raggedright\arraybackslash}p{6.2cm}l@{}}
\toprule
Step & What to do & Evidence in this study & Grade \\
\midrule
1. Define the target & Fix the task, sites, scanners and a worst-site acceptance level before training. & Canine ERM: mean 0.785, worst scanner 0.534; the worst site differed between methods (Sec.~\ref{sec:results-domains}). & strong \\
2. Collect sites & Split by site; size the number of sites from a pilot estimate of $\sigma_\Delta$. & Median ICC 0.72--0.95; a 0.05 gain needs 7--53 sites (tumor) or 3 (mitosis); extra seeds help little (Sec.~\ref{sec:results-variance}, \ref{sec:results-contrasts}). & strong \\
3. Train once, strongly & Fine-tune a pathology foundation model with H\&E jitter; tune its own learning rate; check frozen features on the task; give DG objectives low priority. & Jitter better than ERM in 18/20 domains; DINO FT + jitter best mean worst-site AUROC (0.895); frozen probes 18th--21st on mitosis; BT with inherited learning rate 0.715 vs 0.891 on MIDOG++; no DG objective in the top five (Sec.~\ref{sec:results-ranks}). & moderate \\
4. Adapt without labels & Optionally re-estimate BN statistics on shuffled site images, or normalize stain per slide; never per small tile. & BN-adapt $+0.120$ on canine (5/5 scanners), $+0.015$ and $+0.022$ on Camelyon17 and MIDOG~2021, $-0.004$ on MIDOG++; per-tile Macenko 0.557 vs 0.976 per slide at one center (Suppl.~\ref{S-app:macenko}). & mixed \\
5. Select and validate & Select on training-site data; hold out whole sites with $\ge$3 seeds; report every site, the worst site, seeds, ECE and intervals. & Selection rules agree at the standard tier ($\tau_b\ge0.83$) but not at the compact tier (down to 0.32); ECE 0.027--0.186 on Camelyon17 (Sec.~\ref{sec:results-sensitivity}). & strong \\
6. Accept each new site & Label a small random sample; accept if the site acceptance test gives $P(\text{AUROC}\ge T)\ge0.95$; plan $m$ with its planner; monitor afterwards. & Development results alone: 8.6\% false acceptance at $T=0.80$. Test: false acceptance $\le$2.1\%; at $T=0.80$, 84\% power with 100 labels, which local labels alone need 200 for; worst subgroup 8.7\% and 5.7\% with 25 and 50 labels, 3.3\% with 100 (Sec.~\ref{sec:results-sat}). Monitoring not tested. & moderate \\
\bottomrule
\end{tabular*}
\end{table}

\paragraph{What the roadmap does not promise} It does not guarantee that a
model will work at a new site; it raises the chance, measures the risk and
checks each site before use. Its
evidence comes from patch classifiers at 96 to 224~px on four public research
cohorts, not from whole-slide or detection systems in clinical use, and clinical
deployment additionally needs prospective and regulatory validation that this
study does not provide. The method choices in Step~3 were the best on average,
not on every cohort, which is exactly why Steps 5 and 6 cannot be skipped.

\subsection{What method developers should report}
\label{sec:disc-research}

For researchers who propose or compare methods, the results add five reporting
rules to the roadmap. Each one changed a conclusion in this study.
\begin{enumerate}
  \item \emph{Report every domain and repeat runs.} A single official split is
    one sample from a wide distribution, and the seeds of one method spanned
    up to 0.48 AUROC at a single domain.
  \item \emph{Report intervals and power.} State $\sigma_\Delta$ and the
    number of domains a claimed gain would need, treat a large non-significant
    gain as unresolved, and account for the training data that folds share
    \citep{nadeau2003inference,bates2023cv}.
  \item \emph{Include controls and tune every method.} On MIDOG++,
    compute-matched ERM was the best compact-tier arm and well ahead of
    SimCLR, which passes the same number of images through the encoder; on
    the canine cohort it was the worst. Without the control, neither result
    could be attributed correctly.
  \item \emph{Report all model-selection rules, and separate transductive
    methods.} Methods that use unlabeled test-site images were among the
    strongest and should be compared with each other.
  \item \emph{Do not transfer rankings across tasks, and release per-run
    records.} Cohorts that shared only the shift type ranked methods almost
    independently at the standard tier, and stored records let others add
    rules and metrics without retraining.
\end{enumerate}

%% Case study (withdrawn v1 synergy claim) and the NCT-CRC-HE scope result.
%% All numbers are from the authors' earlier (v1) pipeline, recomputed from:
%%   runs/camelyon17_center2/*.json, results/camelyon17_center2/*.csv
%%   $DATA/runs/fold{0..4}/{baseline,sam,ssl,ssl_sam}_seed{0,1,2}.json (ood_test.auroc)
%%   results/nct_crc_he/crc_results_summary.csv, crc_results_raw.csv,
%%   crc_interaction.csv, crc_data_manifest.json
%% v1 arm names: baseline = ERM, ssl = SimCLR, sam = SAM, ssl_sam = SimCLR+SAM.
%% v2 interaction: results/v2/analysis/paper_numbers.json -> simclr_sam_interaction_C (from C/*/per_run_id.csv).

\subsection{Case study: a single-hospital synergy claim that did not survive}
\label{sec:casestudy}

This study began with a claim of our own that the protocol of
Section~\ref{sec:protocol} overturns. The numbers below come from our earlier
(v1) pipeline: a compact CNN on 64~px crops, five epochs and selection on the
OOD-validation hospital, not the v2 protocol.

\paragraph{The original claim} The model was trained on Camelyon17 hospitals 0,
3 and 4, selected on hospital 1 and tested on hospital 2, the official WILDS
split \citep{koh2021wilds}. With one seed per arm, test AUROC was 0.931 for
ERM, 0.895 for SimCLR, 0.826 for SAM and 0.951 for SimCLR+SAM. The interaction
$I=\delta_{\text{SimCLR+SAM}}-\delta_{\text{SimCLR}}-\delta_{\text{SAM}}$, with
$\delta$ the difference from ERM, was $+0.161$ ($\delta_{\text{SimCLR}}=-0.037$,
$\delta_{\text{SAM}}=-0.105$, $\delta_{\text{SimCLR+SAM}}=+0.019$). Each method
alone was worse than ERM and the pair was better, so we read $I$ as a synergy.
Patch-bootstrap intervals were narrow (SAM$-$ERM:
$[-0.108,-0.102]$).

\paragraph{What the design got wrong} The patch bootstrap measures test
sampling, not run-to-run variation. A second ERM seed in the same experiment
scored 0.881, 0.050 below the first and about 15 times the half-width of that
interval. With both available seeds for ERM and SimCLR, $I$ falls to $+0.112$.
As a contrast of four single runs, $I$ carries the run noise of all four arms.

\paragraph{How the protocol detects it} The v1 grid has all five hospitals and
three seeds per cell (fold $t$ tests on hospital $t$). At hospital 2, $I$ is
$-0.019$. Over hospitals 0--4 it is $-0.012$, $-0.012$, $-0.019$, $+0.180$ and
$-0.050$: negative at four of five, with mean $+0.017$ (95\% CI
$[-0.097,+0.132]$, $p=0.70$, Eq.~\ref{eq:mean}). The single positive value
comes from one SAM run at hospital 3 with test AUROC 0.407 (other seeds:
0.908 and 0.864). At hospital 2, the seed SDs of ERM, SAM,
SimCLR and SimCLR+SAM (0.105, 0.065, 0.048 and 0.002) give a single-run SD of $I$ of about 0.13
if run noise is independent across arms, so $+0.161$ is about 1.2 SDs from
zero. Seed replication thus shows that one run cannot separate the effect from
noise, and domain completeness that its sign depends on the test hospital.
We withdraw the synergy claim. Under the v2 protocol (compact tier, five
seeds per arm, ID-validation selection), the interaction on Camelyon17 is
$-0.003$, $-0.007$, $+0.066$, $+0.007$ and $+0.057$ at hospitals 0--4, with
mean $+0.024$ (95\% CI $[-0.019,+0.067]$, $p=0.20$). On the other cohorts it is
$-0.007$ (canine), $-0.013$ (MIDOG~2021) and $+0.003$ (MIDOG++). Under v2, the
official test hospital shows the largest positive value ($+0.066$). A
single-hospital study on that split would again report a synergy that the full
design does not support. The small negative value on MIDOG~2021 (95\% CI
$[-0.021,-0.005]$, unadjusted $p=0.02$, $K=3$) is the only interval that
excludes zero.

\paragraph{Validity of the v1 records} The v1 defects documented in this
revision (MixStyle active at evaluation, Macenko as one fixed color transform,
RotInv identical to ERM, inactive IRM annealing, an empty GroupDRO group on
MIDOG~2021 and the canine cohort, TENT on single-class test batches) do not
touch these four arms. One deviation does apply: v1 SAM updated batch-norm running statistics on the
perturbed pass instead of the unperturbed one, the reverse of the standard
implementation, and v2 corrects it. It affects SAM and SimCLR+SAM in both
v1 experiments, by an amount we have not measured. The single-hospital
experiment also used an earlier script and a different 30,000-patch training
sample than the grid, so its values are not interchangeable with fold~2's.

\subsection{A single-institution cohort cannot support domain claims (NCT-CRC-HE)}
\label{sec:nct}

We also ran the four v1 arms on NCT-CRC-HE-100K \citep{kather2018nct}, with
two seeds each. The task was tumor epithelium versus all other tissue except
background, on the central $64\times64$~px of each tile. Training (20,000
tiles), ID validation (4,000) and the SimCLR pool (10,226 unlabeled tiles)
came from the color-normalized release. There were two test conditions: the
same release without color normalization (4,000 tiles), and CRC-VAL-HE-7K
\citep{kather2019predicting}, which comes from 50 other patients of the same
tissue bank (class-balanced to 2,466 tiles).

\paragraph{$K=1$ by construction} The tiles have no site labels, so no domain
can be held out. The cohort is one domain, and model selection had to use ID
validation. Neither test condition changes the institution.

\paragraph{Result} All arms reach ID-validation AUROC of 0.963--0.972 (means
over two seeds). Without color normalization, ERM scores 0.791, SAM 0.726,
SimCLR+SAM 0.724 and SimCLR 0.702 ($-0.089$ relative to ERM). On CRC-VAL-HE-7K
the four arms span only 0.017 (0.927--0.944). ERM has the lowest ECE in both
test conditions (0.157 against 0.206--0.235, and 0.027 against 0.032--0.037).
Note that SimCLR's pool saw only normalized tissue, so it never saw the
un-normalized appearance it was tested on.

The arithmetic of Section~\ref{sec:casestudy} recurs. Without normalization,
both single methods are worse than ERM, and $I=+0.087$. A tile bootstrap of
the two-seed ensembles gives $+0.078$ $[0.070,0.085]$. That interval reflects
tile sampling only, while SimCLR's two seeds differ by 0.102 in the same
condition (0.752 and 0.651).

\paragraph{Why this is a scope statement} With $K=1$ there are no
leave-one-domain-out folds and no between-domain variance $\sigma_D$ or ICC
(Eq.~\ref{eq:oneway}). Eq.~\ref{eq:mean} also needs $K\ge2$. The result shows
that, on one tissue source, SimCLR pre-trained on normalized tiles did not
help under an unseen stain appearance or on new patients. It says nothing
about institutional shift either way. We report it to set the scope of our conclusions. The SAM
batch-norm caveat above applies here too.


\subsection{Limitations}
\label{sec:limitations}

\paragraph{Four cohorts, one per cell} The task-by-shift analysis has one
cohort per cell. Task, species, tissue, label definition and training-set size
vary together between cells, and the cohort-level decomposition is
descriptive. The two mitosis cohorts share MIDOG~2021's 150 images. More
cohorts per cell would be needed to separate task from shift type
statistically.

\paragraph{Few domains} With $K=3$ to 7, the ICC and the paired contrasts have
wide intervals, and fold dependence can only be handled as a sensitivity
analysis. MIDOG~2021 has no OOD-validation domain.

\paragraph{Model scale and input size} The standard tier uses ResNet-50 at
96~px and ViT-S/16 at 224~px. Patch classification at this size is not
whole-slide inference or mitotic-figure detection, and larger inputs or
detection pipelines may rank methods differently. The two tiers already
disagree in part (Section~\ref{sec:results-ranks}).

\paragraph{Foundation models} HuggingFace was not reachable from our compute
environment, so recent large foundation models such as UNI, Virchow or
Prov-GigaPath \citep{chen2024uni,vorontsov2024virchow,xu2024gigapath} were not
evaluated. The Lunit models are trained on TCGA only and are licensed for
non-commercial research.

\paragraph{Tuning} Hyperparameters were tuned with one seed per fold, and the
compact tier uses a fixed learning rate. Some selected values lie at a grid
edge (Supplementary Section~\ref{S-app:tuning}). A larger search could change individual
ranks.

\paragraph{Transductive stain normalization} Per-slide Macenko normalization
estimates each test slide's stain from that slide's benchmark patches, which
were sampled from annotated regions. The effect of this tile set on the
estimate was small (Supplementary Section~\ref{S-app:macenko}), but a label-independent tissue
sample would match practice better.

\paragraph{Secondary cohorts} PatchCamelyon overlaps with Camelyon17 in its
hospitals, and NCT-CRC-HE has no site labels, so neither tests institutional
shift. The case studies use numbers from our earlier pipeline, which differs
from the v2 protocol.

\paragraph{Site acceptance test} The test assumes that a new site is
exchangeable with the development sites, so a site unlike all of them must be
caught by the local labels; in the hardest subgroup it needed 100 labels to
keep false acceptance below 5\%. It was validated on labeled patches drawn at
random from each site's test patches, whereas real labels come in clusters from
a few slides, which widens the local uncertainty. It addresses patch-level
AUROC, not a clinical endpoint, and its label planner was slightly optimistic
at the standard tier.

\paragraph{Metric} AUROC is the primary metric. It ignores the decision
threshold, which matters in practice, and threshold-dependent metrics may
rank methods differently. They are stored in every record.
